\documentclass{article}

\usepackage[utf8]{inputenc}
\usepackage[T1]{fontenc}
\usepackage{xcolor}
\usepackage{pifont}
\usepackage{amssymb}
\usepackage{changes}
\newcommand{\done}{\textcolor{green!60!black}{\ding{51}}}
\newcommand{\pending}{$\square$}
\newcommand{\wip}{\textcolor{orange}{$\triangleright$}}
\newcommand{\wontdo}{\textcolor{red}{\ding{55}}}

\usepackage{longtable}

\title{Reconstruction Plan for Honours Thesis}

\begin{document}

\maketitle

\section*{Goal}

Reconstruct a complete and compilable \LaTeX{} source tree for the honours thesis from the
final PDF. The objective is preservation and reproducibility rather than exact page-by-page
reproduction.

\section*{Success Criteria}

\begin{description}
\item[Gold:] Source compiles and reproduces all text, equations, figures, tables,
citations and appendices.

\item[Silver:] Scientific content is fully reproduced but formatting differs from
the original PDF.

\item[Bronze:] PDF preserved, structure documented, and partial source recovered.
\end{description}

Target level for this thesis: \textbf{Gold}.

\section{\done Recreate Structure}

Create a minimal compiling thesis skeleton based on the Table of Contents.

\begin{verbatim}
honours-thesis/
|----- thesis.tex
|----- chapters/
|   |----- intro.tex
|   |----- formulation.tex
|   |----- data.tex
|   |----- modelling.tex
|   |----- results.tex
|   |----- conclusions.tex
|   |----- appendixA.tex
|   |----- appendixB.tex
|   |----- appendixC.tex
|----- figures/
|----- thesis.bib
|----- thesis.cls
\end{verbatim}

\subsection{Status}
\begin{itemize}
\item[\done] Main file created
\item[\done] Chapter files created
\item[\done] Appendices created
\item[\done] Add sections, subsections etc. to files
\item[\done] PDF compiles with book class
\item[\done] Find custom cls file
\item[\done] PDF compiles with custom cls file
\end{itemize}


\section{\done Reconstruct Bibliography}

Extract all bibliography entries from the PDF and recreate a valid BibTeX database.

\subsection{Tasks}
\begin{itemize}
\item[\done] Create \texttt{thesis.bib}
\item[\done] Add placeholder entries for all 69 references
\item[\done] Add bib entries based on ADS searches
\item[\done] Confirm unambiguous ADS/arXiv entries
\item[\done] Confirm entries which need human verification
\item[\done] Recover any missing entries manually
\end{itemize}

\subsection{Bibliography Recovery Status}

\small

\begin{longtable}{|r|p{8cm}|c|}
\hline
Ref & Citation & Status \\
\hline
\endhead
\hline
\endfoot
1 & Albrecht et al. (2006), Report of the Dark Energy Task Force & \done \\
2 & Amanullah et al. (2010), Union2 Compilation & \done \\
3 & Barnes, Francis et al. (2009), Evolving Dark Energy & \done \\
4 & Blandford (2010), New Worlds, New Horizons & \done \\
5 & Blondin et al. (2008), Time Dilation in Type Ia SNe & \done \\
6 & Brewer \& Francis (2009), Entropic Priors & \done \\
7 & Carroll (2001), Cosmological Constant & \done \\
8 & Carroll (2007), Lecture Notes on General Relativity & \done \\
9 & Carroll, De Felice et al. (2005) & \done \\
10 & Carroll \& Kaplinghat (2002) & \done \\
11 & Christensen et al. (2001), Bayesian CMB Methods & \done \\
12 & Copeland, Sami et al. (2006), Dynamics of Dark Energy & \done \\
13 & de Putter \& Linder (2008) & \done \\
14 & do Carmo (1976), Differential Geometry & \done \\
15 & do Carmo (1992), Riemannian Geometry & \done \\
16 & Dunkley et al. (2005), Fast and Reliable MCMC & \done \\
17 & Dunkley et al. (2009), WMAP Five-Year Results & \done \\
18 & Dunlop et al. (1996), 3.5 Gyr Galaxy at z=1.55 & \done \\
19 & Efstathiou (2008), Limitations of Bayesian Evidence & \done \\
20 & Einstein (1917), Cosmological Considerations & \done \\
21 & Einstein (1918), Principle of General Relativity & \done \\
22 & Eisenstein et al. (2005), BAO Detection & \done \\
23 & Francis (2008), PhD Thesis & \done \\
24 & Friedmann (1922), Curvature of Space & \done \\
25 & Gregory (2005), Bayesian Logical Data Analysis & \done \\
26 & Guy et al. (2005), SALT & \done \\
27 & Guy et al. (2007), SALT2 & \done \\
28 & Guy et al. (2010), SNLS 3-Year Sample & \done \\
29 & Hamann \& Ferland (1993) & \done \\
30 & Harrison (1993), Cosmology & \done \\
31 & Hartle (2003), Gravity & \done \\
32 & Hasinger et al. (2002) & \done \\
33 & Hicken et al. (2009), Constitution Set & \done \\
34 & Hobson, Efstathiou \& Lasenby (2006) & \done \\
35 & Hogg (2000), Distance Measures in Cosmology & \done \\
36 & Hogg et al. (2002), K Correction & \done \\
37 & Ivers (2010), Mathematical Computing Notes & \done \\
38 & Jha, Riess et al. (2007), MLCS2k2 & \done \\
39 & Komatsu et al. (2010), WMAP Seven-Year Results & \done \\
40 & Kowalski et al. (2008), Union Compilation & \done \\
41 & Krauss (1997) & \done \\
42 & Kwan, Francis et al. (2009) & \done \\
43 & Lahav \& Liddle (2010), Cosmological Parameters & \done \\
44 & Lampeitl et al. (2010) & \done \\
45 & Leibundgut (2008), Supernovae and Cosmology & \done \\
46 & Leibundgut (2009), Supernova Cosmology 2010 & \done \\
47 & Lewis \& Bridle (2002), Cosmological Parameters from CMB & \done \\
48 & Lidman (2010), Personal Communication & \done \\
49 & Lidman (2010), Improving Type Ia SNe & \done \\
50 & Linder (1988), Cosmological Tests & \done \\
51 & Linder (1988), Light Propagation & \done \\
52 & Linder (2009), Like vs Like & \done \\
53 & Linder \& Scherrer (2008), Aetherizing Lambda & \done \\
54 & Maeda et al. (2010), Asymmetric SN Explosions & \done \\
55 & Mazzali et al. (2007), Common SN Ia Mechanism & \done \\
56 & Olive (2010), Dark Energy and Dark Matter & \done \\
57 & Percival et al. (2007), BAO Scale Measurement & \done \\
58 & Rees (1998), Universe at z > 5 & \done \\
59 & Reid (1998), Globular Clusters and Age of Galaxy & \done \\
60 & Riess et al. (2009), Redetermination of Hubble Constant & \done \\
61 & Riess et al. (2007), HST Type Ia Supernovae & \done \\
62 & Riley, Hobson \& Bence (2002), Mathematical Methods & \done \\
63 & Rubin, Linder et al. (2008) & \done \\
64 & Sawicki \& Carroll (2005) & \done \\
65 & Silk (1989), The Big Bang & \done \\
66 & Spergel et al. (2007), WMAP Three-Year Results & \done \\
67 & van Osten (2006), Dark Energy and Cosmic Horizons & \done \\
68 & Wagner et al. (2008), BAO Constraints & \done \\
69 & Weinberg (2009), Computing the Bayesian Factor & \done \\
\hline
\end{longtable}

\todo{For the theses (van Osten and Barnes) I have tried searching the USyd thesis database but the recor is patchy.  They don't appear to be on arxiv either.}

\section{\done Recover Figures}

Before reconstructing text, create placeholders for every figure.

\subsection{Tasks}
\begin{itemize}
\item[\done] Inventory all figures
\item[\done] Create all figure environments and captions.
\item[\done] Extract figures from PDF where necessary
\item[\pending] Find original figure files from references if available
\item[\pending] Reproduce figures from code if necessary
\item[\done] Create figure labels
\item[\done] Insert figures into environments
\item[\done] Update captions
\end{itemize}

\subsection{Figure Recovery Status}

\begin{longtable}{|l|p{4.5cm}|c|c|c|c|p{3cm}|}
\hline
Figure &
Description &
Cropped? &
Inserted? &
Rebuild? &
Recovered? &
Notes \\
\hline
\endhead

1.1 & Stress-energy-momentum tensor, general and perfect fluid cases & \done & \done & \pending & \pending & \\
1.2 & Supernova classes based on spectroscopy & \done & \done & \pending & \pending & Literature figure \\
1.3 & Type Ia supernova formation diagram & \done & \done & \pending & \pending & Literature/web figure \\

3.1 & Redshift-bin distribution for Gold, Union and Constitution samples & \done & \done & \pending & \pending & From Kwan et al. \\

4.1 & Luminosity distance and distance modulus curves compared with Hogg & \done & \done & \pending & \pending &  Contains imported Hogg plots plus own plots \\
4.2 & Example of proper MCMC maximum step size & \done & \done & \pending & \pending &  MATLAB-generated \\
4.3 & Examples of improper MCMC maximum step sizes & \done & \done & \pending & \pending &  MATLAB-generated \\

5.1 & One-dimensional posterior PDFs for parameters & \done & \done & \pending & \pending &  Also enlarged in Appendix A \\
5.2 & Credible regions for $(\Omega_M,w_0)$ & \done & \done & \pending & \pending &  MATLAB-generated \\
5.3 & Credible regions for $(\Omega_M,w_1)$ & \done & \done & \pending & \pending &  MATLAB-generated \\
5.4 & Credible regions for $(\Omega_M,\gamma)$ & \done & \done & \pending & \pending &  MATLAB-generated \\
5.5 & Credible regions for $(w_0,w_1)$ & \done & \done & \pending & \pending &  Includes CMB wall \\
5.6 & Credible regions for $(w_0,\gamma)$ & \done & \done & \pending & \pending &  MATLAB-generated \\
5.7 & Credible regions for $(w_1,\gamma)$ & \done & \done & \pending & \pending &  MATLAB-generated \\
5.8 & Union and Gold credible-region comparison & \done & \done & \pending & \pending &  Literature figure, recoloured \\
5.9 & Confidence limits on $w(z)$ & \done & \done & \pending & \pending &  Literature figure \\
5.10 & Constraints on $(w_0,w_a)$ from SN, CMB and BAO & \done & \done & \pending & \pending &  Literature figure \\
5.11 & SALT vs MLCS2k2 comparison & \done & \done & \pending & \pending &  Literature figure \\

A.1a & Enlarged posterior for $h/h_{72}$ & \done & \done & \pending & \pending &  Appendix A \\
A.1b & Enlarged posterior for $\Omega_M$ & \done & \done & \pending & \pending &  Appendix A \\
A.1c & Enlarged posterior for $w_0$ & \done & \done & \pending & \pending &  Appendix A \\
A.1d & Enlarged posterior for $w_1$ & \done & \done & \pending & \pending &  Appendix A \\
A.1e & Enlarged posterior for $\gamma$ & \done & \done & \pending & \pending &  Appendix A \\

B.1a-d & Distance-modulus curve plots for Constitution set & \done & \done & \pending & \pending &  Four panels \\
B.2a-d & Distance-modulus curve plots for Constitution 0.2x set & \done & \done & \pending & \pending &  Four panels \\
B.3a-d & Distance-modulus curve plots for Gold set & \done & \done & \pending & \pending &  Four panels \\
B.4a-d & Distance-modulus curve plots for Gold 0.2x set & \done & \done & \pending & \pending &  Four panels \\
B.5a-d & Distance-modulus curve plots for Union set & \done & \done & \pending & \pending &  Four panels \\
B.6a-d & Distance-modulus curve plots for Union 0.2x set & \done & \done & \pending & \pending &  Four panels \\
B.7a-d & Distance-modulus curve plots for SNAP set & \done & \done & \pending & \pending &  Four panels \\
B.8a-d & Distance-modulus curve plots for SNAP 0.2x set & \done & \done & \pending & \pending &  Four panels \\

C.1 & Sample Type Ia supernova light curves & \done & \done & \pending & \pending &  Literature/slide figure \\
C.2 & SALT and MLCS2k2 credible-region comparison & \done & \done & \pending & \pending &  Literature figure \\
C.3 & SN brightness vs host-galaxy star formation rate & \done & \done & \pending & \pending &  Slide figure \\
C.4 & SN brightness vs host-galaxy metallicity & \done  & \done & \pending & \pending &  Slide figure \\
C.5 & Peak absolute magnitude deviation by observing wavelength & \done & \done & \pending & \pending &  Slide figure \\

\hline
\end{longtable}

Captions match the ones in the thesis.
I have added captions to appendices A and B.

\section{\done Recover tables}

\begin{longtable}{|p{1.6cm} | p{5.5cm} c c c p{3cm}|}
\hline
Table &
Description &
Environment? &
Caption/Label? &
Inserted? &
Notes \\
\hline
\endhead
\hline
\endfoot

Table 4.1 &
Programs written for the project &
\done & \done & \done & \\

Table 5.1 &
Best-fit parameters for single-parameter models &
\done & \done & \done & \\

Table 5.2 &
Comparison with original survey results &
\done & \done & \done &  \\

Table 5.3 &
One-parameter deviations from $\Lambda$CDM &
\done & \done & \done & Corrected from original
\end{longtable}

\section{\done Recover Equations}

Reconstruct all displayed equations and numbering.

\subsection{Tasks}
\begin{itemize}
\item[\done] Add correct packages
\item[\done] Add placeholders for displayed equations
\item[\done] Insert displayed equations
\item[\done] All formulae compile
\end{itemize}

\subsection{Equation Recovery Status}

\begin{longtable}{|p{1.8cm}|p{6.2cm}|p{2.2cm}|p{2.2cm}|p{2.5cm}|}
\hline
Equation & Description & Recovered? & Label Added? & Notes \\
\hline
\endfirsthead

\hline
Equation & Description & Recovered? & Label Added? & Notes \\
\hline
\endhead

(1.1) & Einstein field equation & \done & \done & \\
(1.2) & Maximally symmetric spacetime interval & \done & \done & \\
(1.3a) & Robertson form of FLRW metric  & \done & \done & \\
(1.3b) & Walker form of FLRW metric  & \done & \done & \\
(1.4a) & Friedmann acceleration equation  & \done & \done & \\
(1.4b) & Energy conservation equation  & \done & \done & \\
(1.4c) & Friedmann constraint equation  & \done & \done & \\
(1.5) & Einstein equation with cosmological constant  & \done & \done & \\
(1.6) & Alternative Einstein equation with $\Lambda$ term  & \done & \done & \\
(1.7) & Quintessence density equation & \done & \done & \\
(1.8) & Quintessence pressure equation & \done & \done & \\
(1.9) & Combined matter--dark-energy conservation equation & \done & \done & \\
(1.10) & Interacting adiabatic equations & \done & \done & \\
(1.11) & Cosmological redshift relation & \done & \done & \\
(1.12) & Luminosity distance definition & \done & \done & \\
(1.13) & Distance modulus definition & \done & \done & \\
(1.14) & Practical distance modulus expression & \done & \done & \\
(1.15) & Luminosity distance and comoving distance relation & \done & \done & \\
(1.16) & Comoving-distance integral expression & \done & \done & \\

(2.1) & Bayesian product rule & \done & \done & \\
(2.2) & Credible-region probability content & \done & \done & \\
(2.3) & Global likelihood integral & \done & \done & \\
(2.4) & Monte Carlo integral approximation & \done & \done & \\
(2.5) & Change of variable from time to redshift & \done & \done & \\
(2.6) & Cosmological density and critical density & \done & \done & \\
(2.7a) & Scaled matter-density Friedmann equation & \done & \done & \\
(2.7b) & Scaled dark-energy-density Friedmann equation & \done & \done & \\
(2.7c) & Scaled Hubble equation & \done & \done & \\
(2.8) & Luminosity-distance differential equation & \done & \done & \\

(4.1) & Product likelihood and Gaussian individual likelihood & \done & \done & \\

(5.1) & Limiting interacting Friedmann equations as $\gamma \to -1$ & \done & \done & \\

(6.1) & Odds factor and Bayes factor relation & \done & \done & \\
(6.2) & Global likelihood as integral over prior and likelihood & \done & \done & \\

(C.1) & Practical SNe distance modulus relation & \done & \done & \\
(C.2) & SALT light-curve fitter distance modulus relation & \done & \done & \\

\hline
\end{longtable}

\section{\done Recover Text}

Populate chapter files from PDF extraction.

\begin{longtable}{|p{2.2cm}|p{6.5cm}|p{2.2cm}|p{2.2cm}|p{2.5cm}|}
\hline
Section & Title & Text Added? & Reviewed? & Notes \\
\hline
\endfirsthead

\hline
Section & Title & Text Added? & Reviewed? & Notes \\
\hline
\endhead

Abstract & Abstract & \done & \done & \\

1 & Introduction & \done & \done & \\
1.1 & Cosmological background & \done & \done & \\
1.2 & General Relativity & \done & \done & \\
1.3 & Dark energy & \done & \done & \\
1.3.1 & The cosmological constant & \done & \done & \\
1.3.2 & Quintessence models & \done & \done & \\
1.3.3 & Barotropic fluids & \done & \done & \\
1.3.4 & Possible interaction & \done & \done & \\
1.4 & Observations & \done & \done & \\
1.4.1 & Redshift-distance relations & \done & \done & \\
1.4.2 & Supernovae as standard candles & \done & \done & \\
1.4.3 & Relation of observables to parameters & \done & \done & \\
1.5 & Thesis overview & \done & \done & \\

2 & Formulation of the Problem & \done & \done & \\
2.1 & Bayesian inference & \done & \done & \\
2.2 & Determination of model likelihood & \done & \done & \\
2.3 & Markov Chain Monte Carlo methods & \done & \done & \\
2.4 & Reformulation of the Friedmann equations & \done & \done & \\

3 & Data  & \done & \done & \\
3.1 & Observational data  & \done & \done & \\
3.2 & Artificially generated data  & \done & \done & \\

4 & Numerical Modelling  & \done & \done & \\
4.1 & Codes written for the project  & \done & \done & \\
4.1.1 & Numerical solution of the Friedmann equations  & \done & \done & \\
4.1.2 & Likelihood calculations  & \done & \done & \\
4.1.3 & The Metropolis-Hastings algorithm  & \done & \done & \\
4.2 & Generation of the projected WFIRST set  & \done & \done & \\
4.3 & Simulation of the posterior distributions  & \done & \done & \\
4.4 & Credible regions from posteriors  & \done & \done & \\
4.4.1 & Credible regions for individual parameters  & \done & \done & \\
4.4.2 & Credible regions for multiple parameters  & \done & \done & \\

5 & Results  & \done & \done & \\
5.1 & Parameter values from 1-dimensional models  & \done & \done & \\
5.2 & Degeneracies between parameters from 2-dimensional models & \done & \done & \\
5.3 & Improving constraints: combination with CMB and BAO data & \done & \done & \\
5.4 & Implications & \done & \done & \\

6 & Conclusions & \done & \done & \\
6.1 & Further Work  & \done & \done & \\
6.2 & Summary  & \done & \done & \\

A & Large figures from results  & \done & \done & \\
B & Comparison of distance modulus curves to data  & \done & \done & \\
C & Are Type Ia SNe standard candles?  & \done & \done & \\
C.1 & Supernova photometry  & \done & \done & \\
C.2 & The SALT Fitter  & \done & \done & \\
C.3 & The MLCS2k2 Fitter & \done & \done & \\
C.4 & Are the fitting parameters universal?  & \done & \done & \\
C.5 & Future developments  & \done & \done & \\

\hline
\end{longtable}

\section{\done Cross References}

\subsection{Tasks}

For the references within the paper:

\begin{itemize}
\item[\pending] Add packages
\item[\pending] Add labels to chapters and sections.
\item[\pending] Add labels to equations.
\item[\pending] Add labels to figures and tables.
\item[\pending] Add labels to appendices.
\item[\pending] Convert textual references to \texttt{\textbackslash cref} commands.
\item[\pending] Verify all references resolve correctly.
\end{itemize}

For the references to external resources:
\begin{itemize}
\item[\pending] Compile and fix missing citation keys
\item[\pending] Only then polish style
\item[\pending] Compile bibliography successfully
\item[\pending] Find any hyperlinks e.g. URLs in the text, fottnotes and check they work.
\end{itemize}


\subsection*{Cross-Reference Recovery Strategy}

The goal is to recover all internal references in a maintainable manner while avoiding
hard-coded numbering. The recommended approach is to reconstruct labels first and
convert textual references later.

\subsubsection*{Packages}

Load the packages in the following order:

\begin{verbatim}
\usepackage{hyperref}
\usepackage[capitalise,noabbrev]{cleveref}
\end{verbatim}

The \texttt{cleveref} package must be loaded after \texttt{hyperref}.


\subsubsection*{Label Naming Conventions}

Use semantic labels rather than labels based on numbering.

Preferred examples:

\begin{verbatim}
\label{ch:introduction}
\label{sec:general-relativity}
\label{sec:mcmc}

\label{eq:einstein-field}
\label{eq:friedmann-acceleration}

\label{fig:stress-energy-tensor}
\label{fig:sn-classification}

\label{tab:parameter-estimates}
\end{verbatim}

Avoid:

\begin{verbatim}
\label{eq:1.4a}
\label{fig:5.4}
\label{sec:2.3}
\end{verbatim}

These become difficult to maintain if numbering changes.

\subsubsection*{Chapter and Section Labels}

Add labels immediately after chapter and section declarations.

\begin{verbatim}
\chapter{Introduction}
\label{ch:introduction}

\section{General Relativity}
\label{sec:general-relativity}
\end{verbatim}

\subsubsection*{Equation Labels}

Add labels inside equation environments.

\begin{verbatim}
\begin{equation}
R_{\mu\nu} - \frac12 g_{\mu\nu}R
=
\frac{8\pi G}{c^4}T_{\mu\nu}
\label{eq:einstein-field}
\end{equation}
\end{verbatim}

\subsubsection*{Figure Labels}

Always place figure labels after captions.

\begin{verbatim}
\begin{figure}
\centering
\includegraphics{figures/sn-classes}
\caption{Type Ia supernova classification.}
\label{fig:sn-classification}
\end{figure}
\end{verbatim}

\subsubsection*{Temporary Placeholders}

When reconstructing text, unresolved references should be marked explicitly.

Examples:

\begin{verbatim}
[REF: original Fig. 5.4]
[REF: original Eq. (2.7)]
[REF: original Sec. 4.3]
\end{verbatim}

These placeholders are easy to search for later.

\subsection{Verification Checklist}

\begin{longtable}{|p{8cm}|c|}
\hline
Task & Status \\
\hline
\endfirsthead

\hline
Task & Status \\
\hline
\endhead

Load hyperref and cleveref packages & \pending \\
Create chapter labels & \pending \\
Create section labels & \pending \\
Create equation labels & \pending \\
Create figure labels & \pending \\
Create table labels & \pending \\
Create appendix labels & \pending \\
Replace hard-coded equation references & \pending \\
Replace hard-coded figure references & \pending \\
Replace hard-coded section references & \pending \\
Replace hard-coded appendix references & \pending \\
Resolve all undefined references & \pending \\
Verify hyperlinks function correctly & \pending \\
\hline
\end{longtable}

\section{Verification}

Compare reconstructed thesis against original PDF.

\subsection{Checks}

\begin{itemize}
\item[\pending] All chapters present
\item[\pending] All equations present
\item[\pending] All figures present
\item[\pending] All tables present
\item[\pending] Bibliography complete
\item[\pending] No unresolved references
\item[\pending] PDF compiles cleanly
\end{itemize}

\section*{Repository Structure}

Recommended final archival structure:

\begin{verbatim}
honours-thesis/
|----- source/
|----- figures/
|----- bibliography/
|----- thesis.pdf
|----- README.md
|----- LICENSE
|----- reconstruction-log.tex
\end{verbatim}


\end{document}
